\section{Dinâmica e controle da espaçonave com manipulador robótico}

A dinâmica de uma espaçonave equipada com manipulador robótico constitui um tema central nas operações orbitais de proximidade, pois o movimento das juntas não altera apenas a configuração do braço, mas também afeta a atitude da plataforma e, em determinadas situações, o movimento do seu \gls{cdm}. Diferentemente de manipuladores terrestres, cuja base pode ser tratada como aproximadamente fixa para fins de modelagem, o manipulador espacial opera acoplado a uma plataforma livre, sujeita a reações dinâmicas decorrentes das forças e torques internos gerados pelo acionamento das juntas \cite{fonseca2016dynamics,fonseca2019kinematics}.

Nesse cenário, formulações lagrangianas têm sido amplamente empregadas para descrever a dinâmica do sistema acoplado. No trabalho de da Fonseca, Unfried e Lima, por exemplo, a modelagem de uma espaçonave do tipo manipulador durante a fase de aproximação de curta distância para rendezvous e docking/berthing é desenvolvida por meio da formulação lagrangiana com multiplicadores de Lagrange e função dissipativa de Rayleigh \cite{fonseca2016dynamics}. Essa abordagem permite obter não apenas as equações de movimento, mas também as forças e os torques de reação nas juntas, transmitidos ao longo da cadeia cinemática até a plataforma. Os autores mostram que, em microgravidade, esses efeitos perturbam tanto o movimento rotacional da espaçonave quanto a posição do seu \gls{cdm}, o que pode comprometer a precisão requerida nas operações de captura e berthing \cite{fonseca2016dynamics}.

Do ponto de vista do controle, a necessidade de estabilizar a atitude da plataforma enquanto o manipulador executa sua tarefa impõe exigências adicionais ao sistema de guiamento e controle. No estudo de da Fonseca, Unfried e Lima, as simulações computacionais combinam LQR para o movimento orbital linear relativo e PID para a estabilização da atitude durante a operação do manipulador \cite{fonseca2016dynamics}. Esse resultado evidencia que, em sistemas espaciais com manipuladores, o problema de controle não se restringe ao rastreamento de trajetórias articulares, mas envolve também a compensação dos efeitos reativos do braço sobre a plataforma e a preservação das condições necessárias à aproximação final \cite{fonseca2016dynamics}.

A coerência entre cinemática e dinâmica também desempenha papel importante nesse tipo de modelagem. O capítulo de da Fonseca, Pontuschka e Lima destaca que a formulação cinemática de manipuladores do tipo espaçonave deve incluir, além do referencial fixo ao corpo, os referenciais inercial e orbital, de modo a permitir a descrição consistente da atitude, da órbita e do movimento do manipulador dentro de um mesmo problema \cite{fonseca2019kinematics}. Essa observação é relevante porque a formulação dinâmica posterior depende diretamente da definição adequada desses referenciais e das transformações entre eles.

Além das abordagens centradas no sistema acoplado espaçonave--manipulador, a literatura também apresenta contribuições relevantes para o controle de aproximação orbital sob restrições geométricas e operacionais. O trabalho de Fiori et al. propõe uma formulação para modelagem, simulação e controle de rendezvous automatizado nas proximidades de uma estação espacial, considerando a presença de obstáculos e empregando potenciais virtuais e um formalismo geométrico em grupos de Lie \cite{fiori2024modeling}. Embora esse estudo não trate diretamente de um manipulador robótico acoplado, ele contribui para a compreensão de estratégias de guiamento e controle em proximidade orbital, especialmente em cenários nos quais a trajetória e a orientação da espaçonave precisam respeitar restrições posicionais durante a aproximação \cite{fiori2024modeling}.

De forma complementar, trabalhos recentes baseados em aprendizado de máquina têm buscado incorporar informações físicas do sistema a procedimentos de planejamento e resolução de problemas inversos. Os estudos de \cite{santos2022machine,santos2023physics,santos2022uncertainty} indicam que técnicas de machine learning e physics informed machine learning podem melhorar a convergência de algoritmos de otimização e apoiar o cálculo da cinemática inversa e do planejamento de trajetória em robótica espacial. Embora esses trabalhos não tenham como foco principal a formulação completa da dinâmica acoplada da espaçonave com o manipulador, eles sugerem uma tendência de integração entre modelagem física, controle e inteligência computacional em aplicações orbitais.

De modo geral, a literatura mostra que a dinâmica e o controle de espaçonaves com manipuladores robóticos vêm sendo tratados por diferentes perspectivas, que incluem formulações lagrangianas, técnicas clássicas de controle, estratégias de guiamento em proximidade e métodos computacionais orientados por dados. Ainda assim, permanece relevante o desenvolvimento de abordagens que articulem, de maneira consistente, a dinâmica orbital relativa, a atitude da espaçonave perseguidora e o comportamento acoplado do manipulador robótico durante fases críticas de aproximação e interação com o alvo. Essa necessidade está diretamente relacionada ao escopo da presente dissertação.